\chapter{Burping (Eructation)}

\section*{Definition}
The involuntary or voluntary expulsion of gas from the upper gastrointestinal tract through the mouth, caused by swallowed air (aerophagia) or intraluminal gas production, often accompanied by abdominal distension and discomfort.

\section*{What Happens in Your Body}
Burping occurs when swallowed air accumulates in your stomach and is expelled via relaxation of the upper esophageal sphincter. Aerophagia (excessive air swallowing) occurs during eating, drinking, chewing gum, smoking, or anxiety. Supragastric belching involves voluntary air aspiration into the esophagus followed by immediate expulsion without gastric involvement. Gastric fermentation from certain foods or bacterial overgrowth can increase intraluminal gas production.

\section*{Causes}
\begin{itemize}
 \item Aerophagia (eating too quickly, talking while eating, chewing gum, drinking carbonated beverages)
 \item Gastroesophageal reflux disease (GERD)
 \item Functional dyspepsia
 \item Gastroparesis (delayed gastric emptying)
 \item Helicobacter pylori infection
 \item Chronic pancreatitis or pancreatic insufficiency
 \item Celiac disease or lactose intolerance
 \item Small intestinal bacterial overgrowth (SIBO)
 \item Hiatal hernia
 \item Anxiety, stress, or behavioral factors (supragastric belching)
\end{itemize}

\section*{When to See a Doctor}

\begin{itemize}
  \item Symptoms are severe or getting worse over time
 \item Burping accompanied by severe abdominal pain, vomiting, weight loss, or signs of GI bleeding; inability to pass gas with abdominal distension (possible obstruction)
  \item You have other concerning symptoms such as fever, weight loss, or bleeding
\end{itemize}

\section*{Related Symptoms}
\begin{itemize}
 \item Bloating
 \item Abdominal pain
 \item Heartburn
 \item Nausea
 \item Indigestion
\end{itemize}
\textbf{Important:} Your doctor will check for serious underlying conditions when evaluating persistent burping (eructation).

\section*{Self-Care / Home Management}
\begin{itemize}
 \item Rest and avoid activities that make it worse.
 \item Maintain adequate hydration and a balanced diet.
 \item Over-the-counter remedies may provide symptomatic relief (consult a pharmacist or doctor).
 \item Monitor symptoms and seek professional advice if they persist or worsen.
 \item Avoid self-medicating with prescription medications without medical guidance.
\end{itemize}

\section*{How Your Doctor Will Evaluate You}
\begin{itemize}
 \item A thorough history and physical examination are the first steps in evaluation.
 \item Laboratory tests (blood work, urinalysis) may be ordered based on clinical suspicion.
 \item Imaging studies (X-ray, ultrasound, CT, MRI) may be indicated when structural pathology is suspected.
 \item Specialist referral is arranged when the diagnosis remains uncertain or requires specific expertise.
\end{itemize}

\section*{Treatment Options}
\begin{itemize}
 \item Treatment targets the underlying cause identified through diagnostic evaluation.
 \item Supportive care and symptom management are provided while awaiting definitive therapy.
 \item Medications, lifestyle modifications, and physical therapies may be prescribed as appropriate.
 \item Follow-up ensures treatment effectiveness and allows timely adjustments to the care plan.
\end{itemize}

\clearpage